\documentclass[a4paper,12pt,twoside]{memoir}

% Castellano
\usepackage[spanish,es-tabla]{babel}
\selectlanguage{spanish}
\usepackage[utf8]{inputenc}
\usepackage[T1]{fontenc}
\usepackage{lmodern} % scalable font
\usepackage{microtype}
\usepackage{placeins}
\usepackage{listings}
\usepackage{eurosym}
\usepackage{longtable,booktabs}
\usepackage[export]{adjustbox}
\usepackage{dirtree}

\RequirePackage{booktabs}
\RequirePackage[table]{xcolor}
\RequirePackage{xtab}
\RequirePackage{multirow}

% Multi-page tables using
\usepackage{longtable}
\usepackage{tabularx}

% Muestra un directorio en el árbol
\newcommand{\desc}[1]{
	\ldots{} \begin{minipage}[t]{12cm}
		#1{.}
	\end{minipage}
}

% Cell with line break (e.g. \specialcell{Foo\\bar})
\newcommand{\specialcell}[2][c]{%
	\begin{tabular}[#1]{@{}l@{}}#2\end{tabular}}

% Links
\usepackage[colorlinks]{hyperref}
\hypersetup{
	allcolors = {red}
}

% Ecuaciones
\usepackage{amsmath}

% Rutas de fichero / paquete
\newcommand{\ruta}[1]{{\sffamily #1}}

% Párrafos
\nonzeroparskip


% Imagenes
\usepackage{graphicx}
\newcommand{\imagen}[2]{
	\begin{figure}[!h]
		\centering
		\includegraphics[width=0.9\textwidth]{#1}
		\caption{#2}\label{fig:#1}
	\end{figure}
	\FloatBarrier
}
\newcommand{\imagencite}[3]{
	\begin{figure}[!h]
		\centering
		\includegraphics[width=0.9\textwidth]{#1}
		\caption[#2]{#2 \cite{#3}}\label{fig:#1}
	\end{figure}
	\FloatBarrier
}
\newcommand{\imagengrande}[2]{
	\begin{figure}[!h]
		\centering
		\includegraphics[width=1.2\textwidth,center]{#1}
		\caption{#2}\label{fig:#1}
	\end{figure}
\FloatBarrier
}

\newcommand{\imagenflotante}[2]{
	\begin{figure}%[!h]
		\centering
		\includegraphics[width=0.9\textwidth]{#1}
		\caption{#2}\label{fig:#1}
	\end{figure}
}

% Pie de página con enlace
\newcommand{\foothref}[2]{
	\href{#2}{#1}\footnote{#1: \href{#2}{#2}}
}



% El comando \figura nos permite insertar figuras comodamente, y utilizando
% siempre el mismo formato. Los parametros son:
% 1 -> Porcentaje del ancho de página que ocupará la figura (de 0 a 1)
% 2 --> Fichero de la imagen
% 3 --> Texto a pie de imagen
% 4 --> Etiqueta (label) para referencias
% 5 --> Opciones que queramos pasarle al \includegraphics
% 6 --> Opciones de posicionamiento a pasarle a \begin{figure}
\newcommand{\figuraConPosicion}[6]{%
  \setlength{\anchoFloat}{#1\textwidth}%
  \addtolength{\anchoFloat}{-4\fboxsep}%
  \setlength{\anchoFigura}{\anchoFloat}%
  \begin{figure}[#6]
    \begin{center}%
      \Ovalbox{%
        \begin{minipage}{\anchoFloat}%
          \begin{center}%
            \includegraphics[width=\anchoFigura,#5]{#2}%
            \caption{#3}%
            \label{#4}%
          \end{center}%
        \end{minipage}
      }%
    \end{center}%
  \end{figure}%
}

%
% Comando para incluir imágenes en formato apaisado (sin marco).
\newcommand{\figuraApaisadaSinMarco}[5]{%
  \begin{figure}%
    \begin{center}%
    \includegraphics[angle=90,height=#1\textheight,#5]{#2}%
    \caption{#3}%
    \label{#4}%
    \end{center}%
  \end{figure}%
}
% Para las tablas
\newcommand{\otoprule}{\midrule [\heavyrulewidth]}
%
% Nuevo comando para tablas pequeñas (menos de una página).
\newcommand{\tablaSmall}[5]{%
 \begin{table}
  \begin{center}
   \rowcolors {2}{gray!35}{}
   \begin{tabular}{#2}
    \toprule
    #4
    \otoprule
    #5
    \bottomrule
   \end{tabular}
   \caption{#1}
   \label{tabla:#3}
  \end{center}
 \end{table}
}

%
%Para el float H de tablaSmallSinColores
\usepackage{float}

%
% Nuevo comando para tablas pequeñas (menos de una página).
\newcommand{\tablaSmallSinColores}[5]{%
 \begin{table}[H]
  \begin{center}
   \begin{tabular}{#2}
    \toprule
    #4
    \otoprule
    #5
    \bottomrule
   \end{tabular}
   \caption{#1}
   \label{tabla:#3}
  \end{center}
 \end{table}
}

\newcommand{\tablaApaisadaSmall}[5]{%
\begin{landscape}
  \begin{table}
   \begin{center}
    \rowcolors {2}{gray!35}{}
    \begin{tabular}{#2}
     \toprule
     #4
     \otoprule
     #5
     \bottomrule
    \end{tabular}
    \caption{#1}
    \label{tabla:#3}
   \end{center}
  \end{table}
\end{landscape}
}

%
% Nuevo comando para tablas grandes con cabecera y filas alternas coloreadas en gris.
\newcommand{\tabla}[6]{%
  \begin{center}
    \tablefirsthead{
      \toprule
      #5
      \otoprule
    }
    \tablehead{
      \multicolumn{#3}{l}{\small\sl continúa desde la página anterior}\\
      \toprule
      #5
      \otoprule
    }
    \tabletail{
      \hline
      \multicolumn{#3}{r}{\small\sl continúa en la página siguiente}\\
    }
    \tablelasttail{
      \hline
    }
    \bottomcaption{#1}
    \rowcolors {2}{gray!35}{}
    \begin{xtabular}{#2}
      #6
      \bottomrule
    \end{xtabular}
    \label{tabla:#4}
  \end{center}
}

%
% Nuevo comando para tablas grandes con cabecera.
\newcommand{\tablaSinColores}[6]{%
  \begin{center}
    \tablefirsthead{
      \toprule
      #5
      \otoprule
    }
    \tablehead{
      \multicolumn{#3}{l}{\small\sl continúa desde la página anterior}\\
      \toprule
      #5
      \otoprule
    }
    \tabletail{
      \hline
      \multicolumn{#3}{r}{\small\sl continúa en la página siguiente}\\
    }
    \tablelasttail{
      \hline
    }
    \bottomcaption{#1}
    \begin{xtabular}{#2}
      #6
      \bottomrule
    \end{xtabular}
    \label{tabla:#4}
  \end{center}
}

%
% Nuevo comando para tablas grandes sin cabecera.
\newcommand{\tablaSinCabecera}[5]{%
  \begin{center}
    \tablefirsthead{
      \toprule
    }
    \tablehead{
      \multicolumn{#3}{l}{\small\sl continúa desde la página anterior}\\
      \hline
    }
    \tabletail{
      \hline
      \multicolumn{#3}{r}{\small\sl continúa en la página siguiente}\\
    }
    \tablelasttail{
      \hline
    }
    \bottomcaption{#1}
  \begin{xtabular}{#2}
    #5
   \bottomrule
  \end{xtabular}
  \label{tabla:#4}
  \end{center}
}



\definecolor{cgoLight}{HTML}{EEEEEE}
\definecolor{cgoExtralight}{HTML}{FFFFFF}

%
% Nuevo comando para tablas grandes sin cabecera.
\newcommand{\tablaSinCabeceraConBandas}[5]{%
  \begin{center}
    \tablefirsthead{
      \toprule
    }
    \tablehead{
      \multicolumn{#3}{l}{\small\sl continúa desde la página anterior}\\
      \hline
    }
    \tabletail{
      \hline
      \multicolumn{#3}{r}{\small\sl continúa en la página siguiente}\\
    }
    \tablelasttail{
      \hline
    }
    \bottomcaption{#1}
    \rowcolors[]{1}{cgoExtralight}{cgoLight}

  \begin{xtabular}{#2}
    #5
   \bottomrule
  \end{xtabular}
  \label{tabla:#4}
  \end{center}
}


% Plantilla para los casos de uso
\newcommand{\casoDeUso}[9]{
	\begin{longtable}[H]{@{}ll@{}}
		\toprule
		\begin{minipage}[b]{0.23\columnwidth}\raggedright\strut
			\textbf{CU-#1}\strut
		\end{minipage} & \begin{minipage}[b]{0.71\columnwidth}\raggedright\strut
			\textbf{#2}\strut
		\end{minipage}\tabularnewline
		\midrule
		\endhead
		\begin{minipage}[t]{0.23\columnwidth}\raggedright\strut
			\textbf{Versión}\strut
		\end{minipage} & \begin{minipage}[t]{0.71\columnwidth}\raggedright\strut
			1.0\strut
		\end{minipage}\tabularnewline
		\begin{minipage}[t]{0.23\columnwidth}\raggedright\strut
			\textbf{Autor}\strut
		\end{minipage} & \begin{minipage}[t]{0.71\columnwidth}\raggedright\strut
			\nombre\strut
		\end{minipage}\tabularnewline
		\begin{minipage}[t]{0.23\columnwidth}\raggedright\strut
			\textbf{Requisitos asociados}\strut
		\end{minipage} & \begin{minipage}[t]{0.71\columnwidth}\raggedright\strut
			#3\strut
		\end{minipage}\tabularnewline
		\begin{minipage}[t]{0.23\columnwidth}\raggedright\strut
			\textbf{Descripción}\strut
		\end{minipage} & \begin{minipage}[t]{0.71\columnwidth}\raggedright\strut
			#4\strut
		\end{minipage}\tabularnewline
		\begin{minipage}[t]{0.23\columnwidth}\raggedright\strut
			\textbf{Precondición}\strut
		\end{minipage} & \begin{minipage}[t]{0.71\columnwidth}\raggedright\strut
			#5\strut
		\end{minipage}\tabularnewline
		\begin{minipage}[t]{0.23\columnwidth}\raggedright\strut
			\textbf{Acciones}\strut
		\end{minipage} & \begin{minipage}[t]{0.71\columnwidth}\raggedright\strut
			\begin{enumerate}
				\def\labelenumi{\arabic{enumi}.}
				\tightlist
				#6
			\end{enumerate}\strut
		\end{minipage}\tabularnewline
		\begin{minipage}[t]{0.23\columnwidth}\raggedright\strut
			\textbf{Postcondición}\strut
		\end{minipage} & \begin{minipage}[t]{0.71\columnwidth}\raggedright\strut
			#7\strut
		\end{minipage}\tabularnewline
		\begin{minipage}[t]{0.23\columnwidth}\raggedright\strut
			\textbf{Excepciones}\strut
		\end{minipage} & \begin{minipage}[t]{0.71\columnwidth}\raggedright\strut
			\begin{itemize}
				\tightlist
				#8
			\end{itemize}\strut
		\end{minipage}\tabularnewline
		\begin{minipage}[t]{0.23\columnwidth}\raggedright\strut
			\textbf{Importancia}\strut
		\end{minipage} & \begin{minipage}[t]{0.71\columnwidth}\raggedright\strut
			#9\strut
		\end{minipage}\tabularnewline
		\bottomrule
		\caption{CU-#1 #2.}
		\label{CU:#1}
	\end{longtable}
	\newpage
}

\definecolor{mygreen}{rgb}{0,0.6,0}
\definecolor{mygray}{rgb}{0.5,0.5,0.5}
\definecolor{mymauve}{rgb}{0.58,0,0.82}


% Encoding lstlisting
\lstset{literate=
	{á}{{\'a}}1 {é}{{\'e}}1 {í}{{\'i}}1 {ó}{{\'o}}1 {ú}{{\'u}}1
	{Á}{{\'A}}1 {É}{{\'E}}1 {Í}{{\'I}}1 {Ó}{{\'O}}1 {Ú}{{\'U}}1
	{à}{{\`a}}1 {è}{{\`e}}1 {ì}{{\`i}}1 {ò}{{\`o}}1 {ù}{{\`u}}1
	{À}{{\`A}}1 {È}{{\'E}}1 {Ì}{{\`I}}1 {Ò}{{\`O}}1 {Ù}{{\`U}}1
	{ä}{{\"a}}1 {ë}{{\"e}}1 {ï}{{\"i}}1 {ö}{{\"o}}1 {ü}{{\"u}}1
	{Ä}{{\"A}}1 {Ë}{{\"E}}1 {Ï}{{\"I}}1 {Ö}{{\"O}}1 {Ü}{{\"U}}1
	{â}{{\^a}}1 {ê}{{\^e}}1 {î}{{\^i}}1 {ô}{{\^o}}1 {û}{{\^u}}1
	{Â}{{\^A}}1 {Ê}{{\^E}}1 {Î}{{\^I}}1 {Ô}{{\^O}}1 {Û}{{\^U}}1
	{œ}{{\oe}}1 {Œ}{{\OE}}1 {æ}{{\ae}}1 {Æ}{{\AE}}1 {ß}{{\ss}}1
	{ű}{{\H{u}}}1 {Ű}{{\H{U}}}1 {ő}{{\H{o}}}1 {Ő}{{\H{O}}}1
	{ç}{{\c c}}1 {Ç}{{\c C}}1 {ø}{{\o}}1 {å}{{\r a}}1 {Å}{{\r A}}1
	{€}{{\euro}}1 {£}{{\pounds}}1 {«}{{\guillemotleft}}1
	{»}{{\guillemotright}}1 {ñ}{{\~n}}1 {Ñ}{{\~N}}1 {¿}{{?`}}1
}

% Configuración de listing.
% De https://en.wikibooks.org/wiki/LaTeX/Source_Code_Listings
\lstset{ 
	backgroundcolor=\color{white},   % choose the background color; you must add \usepackage{color} or \usepackage{xcolor}; should come as last argument
	basicstyle=\footnotesize,        % the size of the fonts that are used for the code
	breakatwhitespace=false,         % sets if automatic breaks should only happen at whitespace
	breaklines=true,                 % sets automatic line breaking
	captionpos=b,                    % sets the caption-position to bottom
	commentstyle=\color{mygreen},    % comment style
	deletekeywords={...},            % if you want to delete keywords from the given language
	escapeinside={\%*}{*)},          % if you want to add LaTeX within your code
	extendedchars=true,              % lets you use non-ASCII characters; for 8-bits encodings only, does not work with UTF-8
	frame=single,	                 % adds a frame around the code
	keepspaces=true,                 % keeps spaces in text, useful for keeping indentation of code (possibly needs columns=flexible)
	keywordstyle=\color{blue},       % keyword style
	language=Python,                 % the language of the code
	morekeywords={*,...},            % if you want to add more keywords to the set
	numbers=none,                    % where to put the line-numbers; possible values are (none, left, right)
	numbersep=5pt,                   % how far the line-numbers are from the code
	numberstyle=\tiny\color{mygray}, % the style that is used for the line-numbers
	rulecolor=\color{black},         % if not set, the frame-color may be changed on line-breaks within not-black text (e.g. comments (green here))
	showspaces=false,                % show spaces everywhere adding particular underscores; it overrides 'showstringspaces'
	showstringspaces=false,          % underline spaces within strings only
	showtabs=false,                  % show tabs within strings adding particular underscores
	stepnumber=2,                    % the step between two line-numbers. If it's 1, each line will be numbered
	stringstyle=\color{mymauve},     % string literal style
	tabsize=2,	                     % sets default tabsize to 2 spaces
	title=\lstname                   % show the filename of files included with \lstinputlisting; also try caption instead of title
}

\graphicspath{ {./img/} }

% Capítulos
\chapterstyle{bianchi}
\newcommand{\capitulo}[2]{
	\setcounter{chapter}{#1}
	\setcounter{section}{0}
	\chapter*{#2}
	\addcontentsline{toc}{chapter}{#2}
	\markboth{#2}{#2}
}

% Apéndices
\renewcommand{\appendixname}{Apéndice}
\renewcommand*\cftappendixname{\appendixname}

\newcommand{\apendice}[1]{
	%\renewcommand{\thechapter}{A}
	\chapter{#1}
}

\renewcommand*\cftappendixname{\appendixname\ }

% Formato de portada
\makeatletter
\usepackage{xcolor}
\newcommand{\tutor}[1]{\def\@tutor{#1}}
\newcommand{\course}[1]{\def\@course{#1}}
\definecolor{cpardoBox}{HTML}{E6E6FF}
\def\maketitle{
  \null
  \thispagestyle{empty}
  % Cabecera ----------------
\noindent\includegraphics[width=\textwidth]{cabecera}\vspace{1cm}%
  \vfill
  % Título proyecto y escudo informática ----------------
  \colorbox{cpardoBox}{%
    \begin{minipage}{.8\textwidth}
      \vspace{.5cm}\Large
      \begin{center}
      \textbf{TFG del Grado en Ingeniería Informática}\vspace{.6cm}\\
      \textbf{\LARGE\@title{}}
      \end{center}
      \vspace{.2cm}
    \end{minipage}

  }%
  \hfill\begin{minipage}{.20\textwidth}
    \includegraphics[width=\textwidth]{escudoInfor}
  \end{minipage}
  \vfill
  % Datos de alumno, curso y tutores ------------------
  \begin{center}%
  {%
    \noindent\LARGE
    Presentado por \@author{}\\ 
    en Universidad de Burgos --- \@date{}\\
    Tutor: \@tutor{}\\
  }%
  \end{center}%
  \null
  \cleardoublepage
  }
\makeatother

\newcommand{\nombre}{Iván Arjona Alonso} %%% cambio de comando

% Datos de portada
\title{Aplicación Web para la recopilación, tratamiento y visualización de datos públicos}
\author{\nombre}
\tutor{Dr. José Francisco Díez Pastor\\y Dr. Jesús Manuel Maudes Raedo}
\date{\today}

\begin{document}

\maketitle



\cleardoublepage



%%%%%%%%%%%%%%%%%%%%%%%%%%%%%%%%%%%%%%%%%%%%%%%%%%%%%%%%%%%%%%%%%%%%%%%%%%%%%%%%%%%%%%%%



\frontmatter


\clearpage

% Indices
\tableofcontents

\clearpage

\listoffigures

\clearpage

\listoftables

\clearpage

\mainmatter

\appendix

\apendice{Plan de Proyecto Software}

\section{Introducción}

\section{Planificación temporal}

El desarrollo del proyecto se ha llevado a cabo utilizando metodologias ágiles, basandose en la metodología \textit{scrum} con algunas modificaciones (una sóla persona y sin reuniones diarias).

Se aplicó una estrategia de desarrollo incremental, con iteraciones que llamaremos \textit{sprints}.

El resultado de cada iteración es un entregable, sobre el que se discute en la reunión posterior a cada sprint.

Se realizó, en principio, una reunión a la semana con los tutores para exponer las modificaciones realizadas en el sprint anterior y planificar los cambios a realizar en la siguiente iteración. 

Estas tareas están priorizadas por el tiempo estimado de su realización, se puede ver esta estimación en el enlace de cada sprint a sus tareas. 

A continuación se va a realizar un breve resumen de las tareas realizadas en cada una de las iteraciones, así como la duración de cada sprint y el gráfico \textit{burndown} correspondiente.

\subsection{Sprint 0 (16/02/2018 - 02/03/2018)}

Primer sprint del proyecto. En la la reunión de planificación de este primer sprint se discute de forma general de lo que va a tratar el proyecto. 

Las tareas realizadas durante este sprint fueran la creación y configuración del repositorio y sobre todo investigar sobre las tecnologías y herramientas que se podían utilizar.

Se investigaron posibles fuentes de datos para implementar más adelante: INE, sepe, aeat.

La duración fue de dos semanas en lugar de una para poder documentarse sobre todos los aspectos relevantes del proyecto.

\imagen{sprints/sprint0}{Burndown del sprint 0}

\href{https://github.com/IvanArjona/TFG-Datos-publicos/milestone/1?closed=1}{Tareas del sprint 0 en Github}

\subsection{Sprint 1 (03/03/2018 - 08/03/2018)}

Un objetivo de este sprint es investigar alternativas de bases de datos no relacionades que se podrían utilizar. Se ha elegido MongoDB.

El otro objetivo es empezar a implementar prototipos con las fuentes de datos que se habían encontrado en el sprint anterior. Se implementaron prototipos del INE, de la agencia tributaria y del sepe.

\imagen{sprints/sprint1}{Burndown del sprint 1}

\href{https://github.com/IvanArjona/TFG-Datos-publicos/milestone/2?closed=1}{Tareas del sprint 1 en Github}

\subsection{Sprint 2 (09/03/2018 - 15/03/2018)}

El primer objetivo de este sprint es crear la estructura de la página web con Flask. Utilizando un modelo vista-controlador. También se utiliza bootstrap para ahorrar trabajo en el diseño.

Se implementó un prototipo de la carga de datos hacia la base de datos y otro para la descarga de datos desde la base de datos para ser mostrados.

Se hizo una implementación de las fuentes de datos a partir de los prototipos del sprint 1 de modo que se pueda cargar todas las fuentes de manera automática.

\imagen{sprints/sprint2}{Burndown del sprint 2}

\href{https://github.com/IvanArjona/TFG-Datos-publicos/milestone/3?closed=1}{Tareas del sprint 2 en Github}

\subsection{Sprint 3 (16/03/2018 - 22/03/2018)}

En este sprint se sopesaron varias plataformas para hacer el despliegue de la web. De ellas se eligió DigitalOcean y Nanobox.

Se realizó el despliegue utilizando estas plataformas. \href{http://tfg-datos-publicos.nanoapp.io/}{Web desplagada}.

Se investigaron los posibles riesgos de seguridad como inyecciones NoSQL.

Se implementó un formulario para la consulta de datos en la página web. En esta primera aproximación se podía hacer una consulta comparando con una columna de una de las fuentes de datos.

\imagen{sprints/sprint3}{Burndown del sprint 3}

\href{https://github.com/IvanArjona/TFG-Datos-publicos/milestone/4?closed=1}{Tareas del sprint 3 en Github}

\subsection{Sprint 4 (23/03/2018 - 13/04/2018)}

Este sprint coincide con semana santa, por lo que dura una semana más de lo habitual y la carga de trabajo también es mayor.

Se corrigieron errores en los tipos de las fuentes de datos al tratar con números como cadenas. 

Se implementó una forma de descargar las consultas a partir del formulario. 

Se mejoró interfaz gráfica y el formulario de consulta.

Se modificaron las fuentes de datos para corregir errores y añadir el código de municipio a todas ellas para más tarde poder unirlas.

\imagen{sprints/sprint4}{Burndown del sprint 4}

\href{https://github.com/IvanArjona/TFG-Datos-publicos/milestone/5?closed=1}{Tareas del sprint 4 en Github}

\subsection{Sprint 5 (14/04/2018 - 25/04/2018)}

Este sprint se dedicó a empezar a documentar la memoria y corregir algunos errores en la página de consulta como el reenvio de formularios en firefox y la descarga de consultas en json y csv.

También se añadió una descripción a la fuente de datos para explicar de qué se trata cada una en la interfaz web.

\imagen{sprints/sprint5}{Burndown del sprint 5}

\href{https://github.com/IvanArjona/TFG-Datos-publicos/milestone/6?closed=1}{Tareas del sprint 5 en Github}

\subsection{Sprint 6 (26/04/2018 - 02/05/2018)}

El objetivo principal de este sprint fue implementar la posibilidad de juntar varias subconsultas mediante un join. Otra característica implementada es la de avisar al usuario si se ha sobrepasado el límite de columnas especificado, para que pueda filtrar más fino.

También se consideró hacer cambios en el modelo de datos, pero debido a la alta dimensionalidad se ha dejado como estaba en el sprint anterior.

\imagen{sprints/sprint6}{Burndown del sprint 6}

\href{https://github.com/IvanArjona/TFG-Datos-publicos/milestone/7?closed=1}{Tareas del sprint 6 en Github}

\subsection{Sprint 7 (03/05/2018 - 09/05/2018)}

En este sprint se implementó un mapa coroplético para mostrar los valores de cualquier atributo en el mapa agrupando los municipios por su provincia.

Se eliminar los campos duplicados de las consultas que surgían al realizar join de varias subconsultas. Como estos campos repetidos siempre son iguales, se ha optado por eliminarlos en lugar de renombrarlos.

\imagen{sprints/sprint7}{Burndown del sprint 7}

\href{https://github.com/IvanArjona/TFG-Datos-publicos/milestone/8?closed=1}{Tareas del sprint 7 en Github}

\section{Estudio de viabilidad}

\subsection{Viabilidad económica}

\subsubsection{Costes de personal}

El proyecto se ha llevado a cabo por un desarrollador empleado a tiempo parcial durante 4 meses. Se considera un salario neto de \EUR{1000}. (Ver tabla \ref{tab:personal})

\begin{table}[]
	\centering
	\begin{tabular}{@{}ll@{}}
		\toprule
		Concepto & Coste (\euro) \\
		\midrule
		Salario neto & \EUR{1000}  \\
		Retención IRPF (19 \%) & \EUR{360.53} \\
		Seguridad social (28,30 \%) & \EUR{537.00} \\
		\midrule
		Salario bruto & \EUR{1897.53} \\
		\midrule
		Total 4 meses & \EUR{7592.59} \\
		\bottomrule
	\end{tabular}
	\caption{Costes de personal}
	\label{tab:personal}
\end{table}

Los percentajes de cotización a la seguridad social se han calculado a partir del régimen general para 2018 como horas comunes (23.6\% empresa y 4.7\% trabajadores) \cite{misc:cotizacion}.

\subsubsection{Costes de material}

En esta sección se incluyen los costes de software y hardware.

Como hardware se incluye el coste del equipo que se ha utilizado para desarrollar la aplicación (un único pago de \EUR{1000}) y el servidor web (\EUR{5} mensual). (Ver tabla \ref{tab:costeshardware})

\begin{table}[H]
	\centering
	\begin{tabular}{@{}ll@{}}
		\toprule
		Concepto & Coste (\euro) \\
		\midrule
		Ordenador portátil & \EUR{1000}  \\
		Servidor web (mensual) & \EUR{5} \\
		\midrule
		Total 4 meses & \EUR{1020} \\
		\bottomrule
	\end{tabular}
	\caption{Costes de hardware}
	\label{tab:costeshardware}
\end{table}

Como software incluimos la licencia de Windows como único pago y las licencias de Github developer, Codacy Pro y PyCharm como pago mensual. Gitkraken Pro supone un coste de \EUR{41} anual, como no hay pago mensual, se considerará el pago completo. (Ver tabla \ref{tab:costessoftware})

\begin{table}[H]
	\centering
	\begin{tabular}{@{}ll@{}}
		\toprule
		Concepto & Coste (\euro) \\
		\midrule
		Windows 10 Home  & \EUR{145} \\
		GitHub Developer (mensual) & \EUR{7} \\
		Codacy Pro Plan (mensual) & \EUR{15} \\
		PyCharm (mensual) & \EUR{8.90} \\
		GitKraken Pro (anual) & \EUR{41} \\
		\midrule
		Total 4 meses & \EUR{309.60} \\
		\bottomrule
	\end{tabular}
	\caption{Costes de software}
	\label{tab:costessowftware}
\end{table}

\subsubsection{Costes totales}

Tendremos un coste total de \EUR{8922.19}. (Ver tabla \ref{tab:costestotal})

\begin{table}[H]
	\centering
	\begin{tabular}{@{}ll@{}}
		\toprule
		Concepto & Coste (\euro) \\
		\midrule
		Personal  & \EUR{7592.59} \\
		Hardware & \EUR{1020} \\
		Software & \EUR{309.60} \\
		\midrule
		Total & \EUR{8922.19} \\
		\bottomrule
	\end{tabular}
	\caption{Costes de software}
	\label{tab:costestotal}
\end{table}

\subsection{Viabilidad legal}

Con la ayuda de VersionEye \cite{misc:versioneye} se han listado las dependencias del proyecto, con sus correspondientes licencias (Ver tabla \ref{tab:dependencias}). 

\begin{table}[H]
	\centering
	\begin{tabular}{lll}
		\hline
		Dependencia & Versión & Licencia \\ \hline
		beautifulsoup4 & 4.6.0 & MIT \\
		branca & 0.2.0& MIT \\
		chardet & 3.0.4 & MIT \\
		click & 6.7 & BSD \\
		dominate & 2.3.1 & LGPL-3.0 \\
		Flask & 1.0.2 & BSD \\
		Flask-Bootstrap & 3.3.7.1 & BSD \\
		Flask-PyMongo & 0.5.1 & BSD \\
		Flask-WTF & 0.14.2 & BSD \\
		folium & 0.5 & MIT \\
		itsdangerous & 0.24 & BSD \\
		Jinja2 & 2.10 & BSD \\
		MarkupSafe & 1.0 & BSD \\
		numpy & 1.14.3 & BSD \\
		pandas & 0.22.0 & BSD \\
		pymongo & 3.6.1 & Apache-2.0 \\
		requests & 2.18.4 & Apache-2.0 \\
		six & 1.11.0 & MIT \\
		urllib3 & 1.22 & MIT \\
		WTForms & 2.1 & BSD \\
		xldr & 1.1.0 & BSD \\
	\end{tabular}
	\caption{Dependencias}
	\label{tab:dependencias}
\end{table}

De las dependencias usadas en el proyecto, la licencia más restrictiva es LGPL-3.0 \cite{misc:lgpl}. Esta licencia permite el uso de la librería con total libertad, por lo que no nos restringe en la licencia que tenemos que utilizar.

Se ha decidido utilizar la licencia GNU-3.0 \cite{misc:gnu3}. Con esta licencia se permite el uso, distribución y modificación del proyecto, siempre que se mantenga la misma licencia y se acredite al autor original.

\apendice{Especificación de Requisitos}

\section{Introducción}

En este apéndice se describirán los objetivos de la aplicación y se detallarán tanto los requisitos funcionales como los no funcionales.

\section{Objetivos generales}

\begin{itemize}
	\item Integrar varias fuentes de datos públicas en una única base de datos. Estos datos son datos de carácter sociológicos, económicos y demográfico a nivel municipal en España.
	\item Permitir añadir nuevos conjunto de datos de forma sencilla.
	\item Desarrollar un algoritmo para la carga de varias fuentes de datos en una base de datos de manera automatizada.
	\item Desarrollo de una aplicación web que permita la consulta de los datos de manera sencilla y visual.
	\item Facilitar la interpretación de los datos utilizando un mapas coropléticos interactivo.
	\item Desplegar la aplicación web en un servidor de forma que sea fácil de actualizar cada vez que se realice un cambio. Además de funcionar en un entorno local.
\end{itemize}

\section{Catalogo de requisitos}

\subsection{Requisitos funcionales}

\begin{itemize}
	\item \textbf{RF-1 Cargar datos}: El administrador de datos debe poder cargar y actualizar los datos desde sus respectivas fuentes de forma automatizada.
	\item \textbf{RF-2 Consulta}: Los usuarios deben poder consultar información de las fuentes de datos.
	\begin{itemize}
		\item \textbf{RF-2.1}: Podrán realizarse varias subconsultas al mismo tiempo.
		\item \textbf{RF-2.2}: Se podrá seleccionar las columnas resultantes a mostrar.
		\item \textbf{RF-2.3}: Se podrá filtrar según los campos de una de las fuentes de datos.
	\end{itemize}
	\item \textbf{RF-3 Visualizar datos}: Los usuarios deben poder visualizar los datos de una consulta en forma de mapa coroplético.
	\begin{itemize}
		\item \textbf{RF-3.1}: Los datos podrán elegirse de cualquier columna numérica.
		\item \textbf{RF-3.2}: Se podrán mostrar los datos en un mapa a nivel municipal o provincial.
		\item \textbf{RF-3.3}: Los datos de una columna se agregarán utilizando varios métodos (media, suma y cuenta).
	\end{itemize}
	\item \textbf{RF-4 Exportar}: Los usuarios deben poder exportar datos en varios formatos.
	\item \textbf{RF-5 Juntar csv}: Los usuarios deben poder juntar varios archivos csv.
		\begin{itemize}
		\item \textbf{RF-5.1}: Juntar los archivos mediante una columna común.
		\item \textbf{RF-5.2}: Utilizando varios tipos de join (\textit{inner}, \textit{outer}, \textit{left}, \textit{right}).
		\item \textbf{RF-5.3}: Seleccionar la ruta donde se exporta el resultado.
	\end{itemize}
\end{itemize}

\subsection{Requisitos no funcionales}

\begin{itemize}
	\item \textbf{RNF-1 Usabilidad}: La aplicación debe ser fácil de usar e intuitiva para el usuario.
	\item \textbf{RNF-2 Rendimiento}: La aplicación debe cargar en un tiempo aceptable.
	\item \textbf{RNF-3 Mantenimibilidad}: La aplicación debe permitir añadir características de forma sencilla.
	\item \textbf{RNF-4 Compatibilidad}: La aplicación debe funciona correctamente en los navegadores modernos más utilizados (Edge, Chrome, Firefox, Opera y Safari).
	\item \textbf{RNF-5 Responsividad}: La aplicación debe funcionar en pantallas de cualquier tamaño y adaptar su interfaz a cada pantalla.
	\item \textbf{RNF-6 Escalabilidad}: La aplicación debe poder aumentar su rendimiento al aumentar recursos hardware.
	\item \textbf{RNF-7 Facilidad de despliegue}: La aplicación debe poder desplegarse en un servidor de forma sencilla.
	\item \textbf{RNF-8 Software libre}: Utilizar software libre siempre que sea posible.
\end{itemize}

\section{Especificación de requisitos}

\subsection{Diagrama de casos de uso}

\subsection{Descripcion de casos de uso}



\apendice{Especificación de diseño}

\section{Introducción}

\section{Diseño de datos}

\section{Diseño procedimental}

\section{Diseño arquitectónico}



\apendice{Documentación técnica de programación}

\section{Introducción}

En este anexo se explica todo lo que tiene que conocer el programador tanto para instalar y ejecutar la aplicación como para poder seguir con el desarrollo.

\section{Estructura de directorios}

A continuación se van a mencionar los directorios con una breve descripción de su contenido para ayudar a comprender la estructura del proyecto. Además de los directorios, también se van a incluir los archivos más importantes. \\

\dirtree{%
	.1 /.
	.2 config/ \desc{Archivos de configuración}.
	.3 config.py \desc{Fichero de configuración de la base de datos}.
	.2 docs/ \desc{Documentación: memoria y anexos}.
	.3 latex/ \desc{Código en Latex para generar la documentación}.
	.4 img/ \desc{Imágenes de la documentación}.
	.4 tex/ \desc{Secciones de la documentación en Latex}.
	.2 fuentes/ \desc{Paquete para la carga de datos a partir de las fuentes. Un fichero por cada fuente de datos}.
	.2 prototipos/ \desc{Notebooks de IPython utilizados en un principio para implementar las fuentes de datos antes de pasarlo a la aplicación. Es más fácil trabajar con notebooks porque se puede ver el resultado en cada paso}.
	.2 test/ \desc{Pruebas de la aplicación}.
	.3 iterfaz/ \desc{Pruebas de interfaz hechas con selenium}.
	.3 unitarios/ \desc{Pruebas unitarias hechas con unittest}.
	.2 web/ \desc{Raíz de la aplicación web hecha con Flask}.
	.3 forms/ \desc{Formularios de Wtforms}.
	.3 geojson/ \desc{Archivos GeoJSON que delimitan los límites geográficos en la representación de mapas}.
	.3 static/ \desc{Ficheros estáticos}.
	.4 css/ \desc{Hojas de estilo CSS}.
	.4 imagenes/ \desc{Imágenes}.
	.4 js/ \desc{Scripts para el comportamiento dinámico en Javascript}.
	.3 templates/ \desc{Plantillas para generar HTML en Jinja2}.
	.2 actualiza-fuentes.py \desc{Script para actualizar las fuentes de datos en la base de datos}.
	.2 boxfile.yml \desc{Archivo de configuración de Nanobox}.
	.2 joincsv.py \desc{Aplicación para juntar ficheros CSV}.
	.2 requirements.txt \desc{Dependencias de la aplicación (paquetes de Python)}.
	.2 run.py \desc{Script para ejecutar la aplicación web}.
}

\section{Manual del programador}

Esta sección tiene como objetivo explicar los puntos más importantes a tener en cuenta por desarrolladores que puedan seguir ampliando este proyecto.

\subsection{Añadir fuentes de datos}

El funcionamiento de esta aplicación se basa en los datos obtenidos de fuentes públicas, por lo que un punto importante es explicar cómo añadir más fuentes a nuestra aplicación.

La lógica de esta parte de la aplicación se encuentra en el paquete \textit{Fuentes}.

Lo primero que tenemos que hacer al añadir una fuente es crear un nuevo fichero con el nombre de la fuente. En este fichero crearemos una clase con el mismo nombre. Esta clase heredará de \textit{Fuente} (situado en \textit{Fuente.py}).

Esta clase tiene que implementar el método \textit{carga()}. La funcionalidad de este método es devolver un dataframe de Pandas a partir de los datos realizando una consulta a la fuente que se está implementando. Esta clase tiene que tener un constructor para parametrizar las diferentes consultas que se puedan realizar a la misma fuente de datos.

Por ejemplo, observamos que en el SEPE, la url para consultar el número de contratos en 2018 es la siguiente:

\begin{lstlisting}
https://sede.sepe.gob.es/es/portaltrabaja/resources/sede/
datos_abiertos/datos/Contratos_por_municipios_2018_csv.csv
\end{lstlisting}
Y la url para consultar los datos del paro en 2018 es la siguiente:

\begin{lstlisting}
https://sede.sepe.gob.es/es/portaltrabaja/resources/sede/
datos_abiertos/datos/Paro_por_municipios_2018_csv.csv
\end{lstlisting}

Observamos que el principio de la url es el mismo, por lo que lo guardaremos en una variable. Podemos observar que lo único que cambia es el nombre del fichero, en el que también aparece el año. De modo que podemos dividirlo como dos parámetros en el constructor: nombre de fichero y año. Además se incluirá la descripción, que depende de cada consulta.

\begin{lstlisting}[caption=Ejemplo de fuente de datos (Sepe)]
class Sepe(Fuente):
	"""
	Fuente de datos para el Servicio Público de Empleo Estatal
	"""

	def __init__(self, url, anios, tabla, descripcion_):
		url_sepe = 'https://sede.sepe.gob.es/es/portaltrabaja/resources/sede/datos_abiertos/datos/'
		self.url = url_sepe + url
		self.anios = anios
		descripcion = descripcion_ + " Servicio Público de Empleo Estatal."
		super().__init__('sepe', tabla, descripcion)

    @staticmethod
	def procesa_datos(url):
		# Procesa los datos
		df = pd.read_csv(url, sep=';', encoding='latin1', header=1)
		return df

	def carga(self):
		dataframes = []
		for anio in self.anios:
			url = self.url.format(anio)
			df = self.procesa_datos(url)
			dataframes.append(df)
		df = pd.concat(dataframes)
		return df
\end{lstlisting}

Heredando de la clase \textit{Sepe} crearemos una subclase por cada una de las consultas que queramos realizar. En esta subclase añadiremos un constructor sin parámetros, en el que se llamará al constructor de su superclase parametrizándolo.

Por ejemplo, se crearía la subclase \textit{SepeContratos}, parametrizando el constructor del padre con el nombre del fichero y los años disponibles. Así quedaría la clase \textit{SepeContratos}:

\begin{lstlisting}[caption=Ejemplo de fuente de datos (SepeContratos)]
class SepeContratos(Sepe):
	"""
	Contratos por sexo y por edades por cada municipio
	"""
	
	def __init__(self):
		url = 'Contratos_por_municipios_{}_csv.csv'
		anios = range(2006, 2019)
		descripcion = 'Contratos por sexo y por edades.'
		super().__init__(url, anios, 'contratos', descripcion)
\end{lstlisting}

Una vez implementadas estas clases hay que añadirlas a la lista de fuentes de datos para que se carguen al ejecutar el script de actualizar fuentes. Esta lista está en el fichero \textit{\_\_init\_\_.py} del paquetes \textit{Fuentes}.

\begin{lstlisting}[caption=Lista con todas las fuentes]
from fuentes.aeat import AeatRenta
from fuentes.municipios import Municipios
from fuentes.ine import InePoblacion
from fuentes.sepe import SepeContratos, SepeEmpleo, SepeParo
from fuentes.mir import MirCongreso


# Lista de todas las fuentes disponibles
fuentes = [
	Municipios,
	AeatRenta,
	InePoblacion,
	SepeContratos,
	SepeEmpleo,
	SepeParo,
	MirCongreso
]
\end{lstlisting}

\section{Instalación y ejecución del proyecto} \label{instalacionprogramador}

\subsection{Instalación}

\subsubsection{MongoDB}

Antes de empezar con la instalación de la aplicación tenemos que instalar la base de datos.
Se ha utilizado una base de datos no relacional MongoDB.

Concretamente se instaló la versión \textit{3.6.4 Community Server} \cite{misc:mongodb} para Windows.

También se ha utilizado \textit{MongoDB Compass} \cite{misc:mongodb} para visualizar el contenido de la base de datos. Esta herramienta es opcional.

\subsubsection{Python}

Esta aplicación se ha desarrollado utilizando la version 3.6.4 \cite{misc:python3}, por lo que se recomienda utilizar esta versión o una posterior. En cualquier caso debe de instalarse Python 3 o superior para evitar incompatibilidades.

Todos los paquetes utilizados en la aplicación están listados en ficheros \textit{requirements.txt} junto con sus correspondientes versiones.

Habrá que instalar todas las dependencias utilizando pip\footnote{Gestor de paquetes de Python \cite{wiki:pip}}:

\begin{lstlisting}
pip install -r requirements.txt
\end{lstlisting}

\subsection{Ejecución}

\subsubsection{Actualización de la base de datos}

Antes de ejecutar la aplicación tendremos que construir la base de datos con el contenido de todas nuestras fuentes por primera vez.

Para ello simplemente hay que ejecutar el fichero \textit{actualiza-fuentes.py}. Puede tardar un rato debido a la gran cantidad de datos que se van a cargar.

Este paso puede repetirse cada vez que se quiera actualizar las fuentes de datos. Puede ser utilizar ejecutarlo una vez al mes para mantener al día las fuentes con datos mensuales.

\begin{lstlisting}
python actualiza-fuentes.py
\end{lstlisting}

\subsubsection{Servidor web}

El servidor Flask ya se ha instalado como un paquete, por lo que no necesita más instalación.
Para ejecutarlo hay un script \textit{run.py} que nos lanza el servidor en el puerto 5000.

Una vez lanzado podremos acceder a la página web desde \href{http://localhost:5000}{localhost:5000}.

\begin{lstlisting}
python run.py
\end{lstlisting}

\section{Despliegue}

Para desplegar la aplicación se ha optado utilizar como servidor en la nube DigitalOcean y Nanobox como microservicio para facilitarnos la instalación de la máquina en la nube y la configuración del servidor.

\subsection{Instalación}

\subsubsection{DigitalOcean}

Primero tendremos que registrarnos en \foothref{DigitalOcean}{https://www.digitalocean.com/} y obtener el \textit{token} de nuestro usuario. Tendremos que poner una tarjeta de crédito de la que nos cobrarán el importe del servidor.

Podría utilizarse otro proveedor de hosting como \textit{Amazon Web services} o \textit{Google Compute}.

\subsubsection{Nanobox}

Después nos registraremos en \foothref{Nanobox}{https://nanobox.io/} y creamos una nueva aplicación utilizando el token que hemos obtenido antes en DigitalOcean.

Ahora elegiremos el plan que queramos contratar. Para este proyecto será suficiente con el plan más basico de \EUR{5}, 1 CPU y 1GB de ram.

Tras crear la aplicación instalamos el \foothref{cliente de nanobox}{https://docs.nanobox.io/install/}. La primera vez que lancemos un comando nos pedirá los datos para iniciar sesión en nuestra cuenta.

En el fichero \textit{nanobox.yml} tenemos la configuración con los componentes que se van a instalar y los comandos que se ejecutan al desplegar la aplicación.

\subsection{Despliegue}


Una vez todo instalado y configurado, cada vez que queramos actualizar la aplicación del servidor nos situamos en la carpeta del proyecto y lanzamos el siguiente comando:

\begin{lstlisting}
nanobox deploy
\end{lstlisting}

Con esto, de forma transparente para nosotros, se crea una máquina virtual con nuestro proyecto en la que se instalan la base de datos, el entorno de ejecución y los paquetes y se ejecuta la aplicación en el servidor.

\section{Pruebas del sistema}

Se ha realizado una batería de pruebas para verificar el correcto funcionamiento de la aplicación.

\subsection{Pruebas unitarias}

Se han realizado pruebas de cada módulo para comprobar su funcionamiento. Estas pruebas están situadas en la carpeta \textit{/test/unitarios} y se han realizado con el framework de pruebas \textit{Unittest} \cite{misc:unittest}.

Estas pruebas se han centrado en la ejecución de todos los métodos utilizados en la consulta (\textit{consulta.py}) y en la representación de mapas \textit{mapas.py}.

Unittest está integrado en la librería estándar de Python, por lo que no hará falta instalar ningún paquete adicional. 

Para ejecutar los test utilizamos el siguiente comando en la consola:

\begin{lstlisting}
python -m unittest discover test.unitarios
\end{lstlisting}

\subsection{Pruebas de interfaz}

Para probar la interfaz web se ha utilizado \textit{Selenium} con la API que proporciona para Python \cite{misc:seleniumpython}, en lugar de grabar las acciones mediante el plugin de Firefox.

Estas pruebas están situadas en la carpeta \textit{/test/interfaz} y se ejecutan de manera similar a las pruebas unitarias, ya que también están hechas con el framework \textit{unittest}.

Para poder ejecutar estas pruebas primero tenemos que instalar el paquete de Selenium:

\begin{lstlisting}
pip install selenium
\end{lstlisting}

Una vez instalado podemos ejecutar las pruebas con el siguiente comando:

\begin{lstlisting}
python -m unittest discover test.interfaz
\end{lstlisting}

\subsection{Todas las pruebas}

También podemos ejecutar todas las pruebas al mismo tiempo:

\begin{lstlisting}
python -m unittest discover
\end{lstlisting}
\apendice{Documentación de usuario}

\section{Introducción}

En este apéndice se explica cuáles son los requisitos que deberá cumplir el usuario para ejecutar la aplicación, como instalarla y cómo utilizarla.

\section{Requisitos de usuarios}

Como se trata de una aplicación web, el único requisito que necesita el usuario es un navegador web que soporte, \textit{Javascript}, \textit{cookies} y hojas de estilo \textit{CSS}. Como nuestra aplicación también utiliza \textit{jQuery}, nuestros navegadores soportados serán los mismos \cite{misc:jquerybrowsers}:

\begin{itemize}
	\item Google Chrome.
	\item Microsoft Edge.
	\item Mozilla Firefox.
	\item Internet Explorer 9 o superior.
	\item Safari para Mac.
	\item Opera.
	\item Navegador de Android 4.0 o superior.
	\item Safari para iOS 7 o superior.
\end{itemize}

En el caso de la aplicación para juntar ficheros CSV, haría falta tener instalado Python 3 y \textit{Pandas 0.23}.

\section{Instalación}

\subsection{Aplicación web}

Como se trata de una aplicación web, los usuarios no tendrán que realizar ninguna instalación. Sólo serían necesario acceder a la aplicación desde la siguiente url:

\href{https://tfg-datos-publicos.nanoapp.io}{https://tfg-datos-publicos.nanoapp.io}

En caso de que se quisiera instalar la aplicación en local habría que seguir los pasos explicados en la documentación para programador (\ref{instalacionprogramador}). 

\subsection{Aplicación para juntar CSV}

La otra herramienta proporcionada en este trabajo es la aplicación para juntar ficheros CSV. En este caso, tampoco haría falta instalación como tal, pero si tener instalado Python 3 y Pandas 0.23. Una forma rápida de instalarlo es usando el instalador de \textit{Anaconda}\footnote{Anaconda: \href{https://www.anaconda.com/download/}{https://www.anaconda.com/download/}} con \textit{Python 3.6}.

Para ejecutar la aplicación:

\begin{lstlisting}
python joincsv.py
\end{lstlisting}

\section{Manual del usuario}

En esta sección se va a explicar al usuario cómo utilizar la aplicación.

\subsection{Consulta} \label{consulta}

Para empezar, nos dirigiremos a la \foothref{página de consulta}{https://tfg-datos-publicos.nanoapp.io/consulta}.

\imagen{manual/consulta}{Página de consulta}

En la primera columna, seleccionaremos la fuente de datos de la que queramos consultar la información y los filtros que queramos. Por ejemplo, si queremos consultar los datos de la renta en Miranda de Ebro, seleccionaremos rellenaríamos lo siguiente:

\imagen{manual/consulta-basica-1}{Filtros de la consulta}

Podemos filtrar por cualquier columna de la fuente, utilizando los siguientes comparadores: igual (=), distinto (!=), menor (<), mayor (>), menor o igual (<=), mayor o igual (>=).

En la columna de la derecha nos mostrará la descripción de la fuente que hayamos seleccionado y podremos seleccionar las columnas que queramos ver.

Por último, elegiremos el número de filas a mostrar, este es el número máximo de filas que se pueden ver en la previsualización de la tabla, al descargar la consulta o mostrar el mapa, aparecerán todos los datos (por defecto 1000).

\imagen{manual/consulta-basica-2}{Seleccionar columnas a mostrar}

Una vez rellenados todos los datos para la consulta, le damos al botón `Consultar' y nos llevará a la visualización.

\subsection{Visualización} \label{visualizacion}

Tras haber realizado la consulta, podremos ver los resultados. Se mostrará una tabla con los datos filtrados.

Podemos elegir el número de filas mostradas por cada página, ordenar los datos por la columna que queramos o buscar dentro de la tabla. 

\imagen{manual/visualizacion}{Visualización de la consulta}

En caso de que hayamos superado el número de filas a mostrar, nos aparecerá un mensaje informándonos de cuántas filas se están mostrando y cuál es el número real de filas.

\imagen{manual/limite}{Límite de filas superado}

\subsection{Exportación de consultas} \label{exportar}

Para exportar una consulta, simplemente presionamos el botón `csv' o `json', dependiendo del formato que nos interese.

Se descargará los datos de la consulta completa. Pueden ser más de los que aparecen en la tabla, si ha superado el número de filas a mostrar que hemos elegido al realizar la consulta.

\imagen{manual/exportar}{Exportar consulta}

\subsection{Consulta de varias fuentes}

Empezamos de nuevo en la \foothref{página de consulta}{https://tfg-datos-publicos.nanoapp.io/consulta}.

Una vez rellenado la primera consulta, como se ha explicado en \ref{consulta}, pulsaremos el botón `+' y nos aparecerá otra pestaña con una nueva consulta.

\imagen{manual/subconsulta-boton}{Botón para añadir una subconsulta}

\imagen{manual/subconsulta-tab}{Pestañas con varias subconsultas}

Las consultas se rellenan independientemente igual que en \ref{consulta}. Podemos cambiar de subconsulta pulsando su pestaña.

Hay que tener en cuenta que al seleccionar las columnas a mostrar, tenemos que seleccionar `Codigo Municipio', que es la columna común entre todas las subconsultas por la que se agruparán.

Una vez rellenados los campos de todas las subconsultas, elegiremos cómo queremos que se juntes. Para ello en el campo \textit{join}, que habrá aparecido al presionar `+' por primera vez, seleccionamos el método por el que queremos que se junten.

Los posibles métodos son: \textit{inner}, \textit{inner join}, \textit{outer join}, \textit{left join}, \textit{right join}.

\imagen{manual/subconsulta-join}{Método para juntar subconsultas}

Cuando hayamos terminado pulsamos `Consultar' y el resultado aparecerá de igual modo que en \ref{visualizacion}.

\subsection{Columna calculada} \label{calculada}

Al hacer una consulta, podemos añadir un campo combinando varias columnas como queramos.

Para ello, utilizaremos el campo `Columna calculada'.

Al empezar a escribir el nombre de una columna existente, se autocompletará con las columnas que contengan esos carácteres. Para que esto funcione, habrá que escribir un espacio antes de cada campo.

\imagen{manual/calculada-autocompletar}{Autocompletar en columnas calculadas}

En este campo podemos utilizar cualquier operador matemático. Incluyendo sumas, restas, divisiones, multiplicaciones, potencias, mínimo (\textit{.min()}), máximo (\textit{.max()}), media (\textit{.mean()}), entre otros.

Por ejemplo, usando la fuente `aeat\_renta', si queremos saber el porcentaje de habitantes que han hecho declaraciones, usaremos la siguiente consulta calculada:

\begin{lstlisting}
Número de declaraciones / Número de habitantes * 100
\end{lstlisting}

\imagen{manual/calculada-resultado}{Visualización de una consulta con una columna calculada}

Es importante destacar que hay que seleccionar todas las columnas que usemos en una consulta calculada en el campo `columnas a mostrar'.

Si la consulta calculada no es válida, nos mostrará un mensaje y no se mostrará, pero sí que se mostrará el resto de la consulta.

\imagen{manual/calculada-error}{Columna calculada no válida}

\subsection{Mapas coropléticos}

Para visualizar los mapas coropléticos, primero tendremos que haber realizado una consulta.

En la página donde se muestra el resultado de la consulta, aparece una serie de desplegable para mostrar los mapas. 

\imagen{manual/mapas-consulta-botones}{Desplegables para visualizar el mapa}

El primer desplegable indica el método por el que agrupar los datos. Posibilidades: \textit{mean}, \textit{sum}, \textit{count}.

El segundo desplegable es el nivel al que se representa el mapa. Puede ser a nivel de provincia o a nivel de municipio.

El tercer desplegable `Visualizar mapa', es la columna que se quiere visualizar, se mostrarán sólo las columnas numéricas. Al pulsarlo nos abrirá el mapa visualizando la columna que pulsemos, con las configuración de los dos desplegables anteriores.

\imagen{manual/mapas-consulta-desplegable}{Desplegable para seleccionar la columna a visualizar}

Es importante que si queremos visualizar el mapa por provincia hayamos seleccionado `Codigo Provincia' en columnas a mostrar, y si se va a representar el mapa a nivel de municipio, se deberá haber seleccionado `Codigo Municipio'. 

En este caso visualizaremos la consulta calculada que habíamos elegido en la sección \ref{calculada}, agrupando por media (\textit{mean}).

Mapa por provincias (Figura \ref{fig:manual/mapa-provincias}).

\imagen{manual/mapa-provincias}{Mapa por provincias de la columna calculada}

Mapa por municipios (Figura \ref{fig:manual/mapa-municipios}).

\imagen{manual/mapa-municipios}{Mapa por municipios de la columna calculada}

En el caso de mostrar el mapa por municipios, esta opción sólo funciona el Firefox, en Chrome no funciona por limitaciones del framework (lineas futuras).

\subsection{Juntar ficheros CSV}

Se ha hecho una herramienta para juntar ficheros CSV. Se ha realizado como una aplicación de escritorio para que los usuarios no tengan que subir sus ficheros a una aplicación web por cuestiones de confianza y privacidad.

Para iniciarla:

\begin{lstlisting}
python joincsv.py
\end{lstlisting}

\imagen{manual/join-inicial}{Aplicación de escritorio para juntar ficheros CSV}

Una vez iniciada, habrá que pulsar el botón `Subir csv' y seleccionar todos los ficheros que quiera juntar. Otra opción es seleccionarlos de uno en uno, pulsando varias veces `Subir csv'.

Podemos usar los ficheros resultantes de exportar como CSV de nuestra aplicación (\ref{exportar}) junto con algún fichero del usuario.

Tras subir los ficheros, nos mostrará cuales son los ficheros mostrados (Figura \ref{fig:manual/join-subidos}).

\imagen{manual/join-subidos}{Selección de columnas al juntar CSV}

Si nos hemos confundido y queremos volver a empezar, pulsamos el boton `Vaciar' y volveremos al principio de la aplicación.

Se mostrará una lista de selección múltiple con las columnas comunes a todos los ficheros. Ahora tendremos que seleccionar cuáles son las columnas por las que queremos hacer join (por defecto se seleccionar `Codigo Municipio', si existe en todos los ficheros) y el tipo de join que queramos utilizar (\textit{inner}, \textit{outer}, \textit{left}, \textit{right}).

Cuando hayamos terminado, pulsamos `Join' y nos aparecerá una ventana para seleccionar la ruta donde guardar el fichero resultante.




\bibliographystyle{plain}
\bibliography{bibliografiaAnexos}

\end{document}
